% =====================================================================
%  SIMULAZIONE DI MATEMATICA - 1SP -- IIS "David Maria Turoldo", Zogno (BG)
%  Classe: exam  |  Compilazione: pdflatex (due passaggi)
%  Soluzioni:        decommentare \printanswers
%  Versione BES/DSA: impostare \versioneBEStrue
%  Struttura:        \iniziofila{A} \begin{esercizi} ... \end{esercizi} \finefila
%                    \iniziofila{B} \begin{esercizi} ... \end{esercizi} \finefila
%  Ogni esercizio:   \begin{esercizio}
	%                      \obiettivo{...}
	%                      \question[punti]  (oppure \question + parts con \part[punti])
	%                      testo ...
	%                    \end{esercizio}
%  Nessuno spazio per le risposte: gli studenti scrivono sul foglio a protocollo.
% =====================================================================
\documentclass[11pt,a4paper,addpoints]{exam}

% ---------- Lingua e codifica ----------
\usepackage[utf8]{inputenc}
\usepackage[T1]{fontenc}
\usepackage[italian]{babel}
\usepackage{lmodern}
\usepackage{microtype}

% ---------- Matematica, grafica, tabelle ----------
\usepackage{amsmath,amssymb}
\usepackage{siunitx}
\sisetup{output-decimal-marker={,}, per-mode=symbol, group-separator={\,}}
\usepackage{tikz}
\usetikzlibrary{arrows.meta,calc,patterns,angles,quotes}
\usepackage{pgfplots}
\pgfplotsset{compat=1.18}
\usepackage{graphicx}
\usepackage{multicol}
\usepackage{array,tabularx,booktabs}
\usepackage{xcolor}

% ---------- Pagina ----------
\usepackage[a4paper,top=2cm,bottom=2.2cm,left=2cm,right=2.2cm,
headheight=30pt,headsep=12pt]{geometry}

% ---------- Opzioni ----------
%\printanswers              % Decommentare per mostrare le soluzioni.
\ifdefined\CORRETTORE\printanswers\fi % Opzione per compilare il correttore.
\newif\ifversioneBES
\versioneBESfalse              % <-- \versioneBEStrue per la versione BES/DSA
\newif\iftabellabonus
\tabellabonusfalse

% ---------- Interlinea ----------
%  \interlinea      : testo normale
%  \interlineaparts : sottopunti a), b), ... (piu' ariosa, separa le equazioni)
\newcommand{\interlinea}{1.05}
\newcommand{\interlineaparts}{1.05}
\ifversioneBES
\usepackage[scaled=1.05]{helvet}
\renewcommand{\familydefault}{\sfdefault}
\renewcommand{\interlinea}{1.5}
\renewcommand{\interlineaparts}{1.7}
\fi
%\linespread{1.5}

% ---------- Dati della prova (da compilare) ----------
\newcommand{\materia}{Matematica}
\newcommand{\classe}{1SP}
\newcommand{\indirizzo}{Liceo Scientifico a indirizzo sportivo}
%   Liceo Scientifico | Liceo Scientifico a indirizzo sportivo | Liceo delle Scienze Umane
\newcommand{\tipoprova}{Simulazione di verifica - Versione sintetica}
\newcommand{\argomento}{Numeri interi e razionali}
\newcommand{\durata}{vedi indicazioni}
\newcommand{\annoscolastico}{2026/2027}
\newcommand{\docente}{Prof. Rubens Longhi}
\newcommand{\fila}{A}
\newcommand{\indicazioni}{Svolgi sul foglio a protocollo. Mostra i passaggi essenziali e riduci i risultati. Nelle divisioni tra monomi considera le lettere al denominatore diverse da zero. \textbf{È consentito l'uso della calcolatrice.}\ifversioneBES\ Sono consentiti gli strumenti compensativi previsti dal PDP.\fi}
% ---------- Intestazione e pie' di pagina ----------
%  Ogni fila ricomincia da pagina 1: la sua prima pagina ha l'intestazione
%  della scuola e la numerazione "Pagina x di y" e' calcolata sulla singola fila.
\pagestyle{headandfoot}
\firstpageheader
{\scriptsize\textsc{IIS ``David Maria Turoldo''}\\[-1pt]
	\scriptsize Via Ronco 11 -- 24019 Zogno (BG)}
{}
{\scriptsize A.S.\ \annoscolastico\\[-1pt]\scriptsize \docente}
\firstpageheadrule
\runningheader{\small\materia\ -- \classe}{\small\tipoprova}{\small Fila \fila}
\runningheadrule
\newcommand{\paginefila}{\pageref{fine@fila@\fila}}
\firstpagefooter{}{\small Pagina \thepage\ di \paginefila}{}
\runningfooter{}{\small Pagina \thepage\ di \paginefila}{}

% ---------- Formato degli esercizi ----------
%  Riga 1:  "1."  in grassetto + obiettivo in grigio ........ (n punti)
%  Riga 2+: testo dell'esercizio
% ---------- Formato degli esercizi ----------
\newcommand{\obiettivocorrente}{}
\newcommand{\obiettivo}[1]{%
	\gdef\obiettivocorrente{#1}}
\newcommand{\stampaobiettivo}{%
	\ifx\obiettivocorrente\empty\else
	{\footnotesize\color{black!55}\textit{Obiettivo: \obiettivocorrente}}%
	\fi}
\qformat{\textbf{\thequestion.}\enspace\stampaobiettivo\hfill\colonnapunti{(\totalpoints\ \points)}\gdef\obiettivocorrente{}}
\bonusqformat{\textbf{\thequestion.}\enspace{\footnotesize\color{black!55}\textit{(facoltativo)}}\enspace\stampaobiettivo\hfill\colonnapunti{(+\totalbonuspoints\ bonus)}\gdef\obiettivocorrente{}}
\renewcommand{\questionlabel}{}
\renewcommand{\partlabel}{\thequestion\thepartno)}
\renewcommand{\subpartlabel}{\thesubpart.}
\renewcommand{\choicelabel}{\thechoice)}
\pointpoints{punto}{punti}
\bonuspointpoints{punto bonus}{punti bonus}
\newcommand{\colonnapunti}[1]{{\small #1}}
\pointformat{\small(\thepoints)}
\bonuspointformat{\small(+\thepoints)}

% Sottopunti con interlinea maggiore (le soluzioni tornano all'interlinea normale)
\renewcommand{\partshook}{\linespread{\interlineaparts}\selectfont\setlength{\itemsep}{3pt}\setlength{\rightmargin}{1cm}}

% Numerazione degli esercizi: ogni fila riparte da 1 (esercizio 1, 2, ...);
% il contatore interno di exam prosegue da una fila all'altra, cosi' punti,
% totali e tabelle restano corretti. Usare l'ambiente "esercizi" al posto di "questions".
\newcounter{offsetfila}
\renewcommand{\thequestion}{\number\numexpr\value{question}-\value{offsetfila}\relax}
\newenvironment{esercizi}
{\begin{questions}\setcounter{question}{\value{offsetfila}}}
	{\end{questions}}

% ---------- Esercizi non spezzati tra due pagine ----------
%  \begin{esercizio} ... \end{esercizio} racchiude \obiettivo + \question + parts.
%  Se non ci sta, l'esercizio passa intero alla pagina successiva.
%  Nella versione con soluzioni l'esercizio puo' spezzarsi (i correttori sono lunghi).
\cancelspace   % le soluzioni nascoste non lasciano spazio ne' punti di interruzione
\newenvironment{esercizio}
{\par\ifprintanswers\else\begin{samepage}\fi}
	{\ifprintanswers\else\end{samepage}\fi\par\vspace{3pt}}

% ---------- Soluzioni e scelte multiple ----------
\renewcommand{\solutiontitle}{\noindent\textbf{Soluzione.}\enspace}
\unframedsolutions
\SolutionEmphasis{\color{blue!60!black}\linespread{\interlinea}\selectfont}
\CorrectChoiceEmphasis{\color{blue!60!black}\bfseries}
\checkboxchar{$\square$}
\checkedchar{$\blacksquare$}

% ---------- Tabella punteggi ----------
\hqword{\textbf{Esercizio}}
\hpword{\textbf{Punti massimi}}
\hsword{\textbf{Punti ottenuti}}
\htword{\textbf{Totale}}
\chqword{\textbf{Esercizio}}
\chpword{\textbf{Punti massimi}}
\chbpword{\textbf{Bonus}}
\chsword{\textbf{Punti ottenuti}}
\chtword{\textbf{Totale}}

% ---------- Macro utili ----------
\newcommand{\R}{\mathbb{R}}
\newcommand{\N}{\mathbb{N}}
\newcommand{\Z}{\mathbb{Z}}
\newcommand{\Q}{\mathbb{Q}}
\newcommand{\dd}{\,\mathrm{d}}

% ---------- Inizio e fine di ogni fila ----------
%  \iniziofila{A} ... \finefila   poi   \iniziofila{B} ... \finefila
%  \iniziofila: nuova pagina 1, intestazione scuola, titolo, nome, indicazioni.
%  \finefila:   tabella punteggi della sola fila, voto, firma.
\newcommand{\iniziofila}[1]{%
	\clearpage
	\renewcommand{\fila}{#1}%
	\setcounter{page}{1}%
	\setcounter{offsetfila}{\value{totalq@fila}}%
	\begin{center}
		{\small\bfseries \tipoprova\ifprintanswers\ -- Soluzioni e criteri di correzione\fi}\\[3pt]
		{\Large \argomento}\\[2pt]
		{\small \indirizzo\ -- Classe \classe\ -- Durata: \durata\ -- \textbf{Fila #1}}
	\end{center}
	\vspace{4pt}
	\noindent
	\begin{tabularx}{\textwidth}{@{}l X l p{3.2cm}@{}}
		\textbf{Nome e Cognome:} & \hrulefill & \textbf{Data:} & \hrulefill \\[8pt]
	\end{tabularx}
	\vspace{2pt}
	\noindent\fbox{\parbox{\dimexpr\textwidth-2\fboxsep-2\fboxrule}{\small\linespread{1.2}\selectfont
			\textbf{Indicazioni.} \indicazioni\ Punteggio totale: \pointsinrange{fila#1} punti.}}
	\vspace{2pt}
	\begingradingrange{fila#1}}
\newcounter{totalq@fila}
\newcommand{\finefila}{%
	\setcounter{totalq@fila}{\value{question}}%
	\endgradingrange{fila\fila}%
	\vfill
	\begin{center}\small
		\iftabellabonus
		\partialcombinedgradetable{fila\fila}[h][questions]%
		\else
		\partialgradetable{fila\fila}[h][questions]%
		\fi
	\end{center}
	\vspace{6pt}
	\noindent{\small\textbf{Voto:} \rule{3cm}{0.4pt}\hfill
		\textbf{Firma del docente:} \rule{5cm}{0.4pt}}%
	\label{fine@fila@\fila}}


% Adattamenti per la simulazione: nessun punteggio, soluzioni attivabili.
\renewcommand{\colonnapunti}[1]{}
\qformat{\textbf{\thequestion.}\enspace\parbox[t]{0.95\linewidth}{\stampaobiettivo}\hfill\gdef\obiettivocorrente{}}
\runningheader{\small\materia\ -- \classe}{\small Simulazione}{\small $\Z$, $\Q$ e applicazioni}
\renewcommand{\iniziofila}[1]{%
\clearpage\renewcommand{\fila}{#1}\setcounter{page}{1}
\begin{center}
{\small\bfseries \tipoprova\ifprintanswers\ -- Soluzioni\fi}\\[3pt]
{\Large \argomento}\\[2pt]
{\small Valore assoluto, M.C.D. e m.c.m., proporzionalità e percentuali}\\[2pt]
{\small \indirizzo\ -- Classe \classe}
\end{center}
\noindent\textbf{Nome e cognome:} \hrulefill\quad\textbf{Data:} \rule{2.3cm}{0.4pt}\par\medskip
\noindent\fbox{\parbox{\dimexpr\textwidth-2\fboxsep-2\fboxrule}{\small
\textbf{Indicazioni.} Svolgi sul foglio a protocollo, mostrando i passaggi e riducendo le frazioni ai minimi termini. Lavora senza calcolatrice. Nei vero/falso, giustifica le affermazioni vere e fornisci un controesempio per quelle false. La raccolta contiene due proposte per obiettivo ed è più ampia della verifica di 60 minuti.}}
\medskip}
\renewcommand{\finefila}{\label{fine@fila@\fila}}
\newcommand{\sezione}[1]{\par\medskip\noindent{\large\bfseries #1}\par\smallskip}
\begin{document}
\linespread{\interlinea}\selectfont
\iniziofila{S}
\ifprintanswers\small\fi
\begin{esercizi}
\begin{esercizio}
\obiettivo{distinguere somme nulle e prodotti nulli}
\question Vero o falso? Giustifica se vero, fornisci un controesempio se falso.
\begin{parts}
\part Se $a,b\in\Z$ e $|a|+|b|=0$, allora $a=0$ e $b=0$.
\begin{solution}
\textbf{Vera.} I moduli sono non negativi: se uno fosse positivo, anche la somma sarebbe positiva. Dunque entrambi i moduli, e quindi entrambi i numeri, sono zero.
\end{solution}
\part Se $a,b\in\Z$ e $|a|\cdot|b|=0$, allora $a=0$ e $b=0$.
\begin{solution}
\textbf{Falsa.} $a=0$, $b=5$ è un controesempio. La legge di annullamento del prodotto richiede che \emph{almeno uno} dei fattori sia zero.
\end{solution}
\end{parts}
\end{esercizio}

\begin{esercizio}
\obiettivo{argomentare su valore assoluto e numeri razionali}
\question Vero o falso? Giustifica se vero, fornisci un controesempio se falso.
\begin{parts}
\part Per ogni $a,b\in\Q$ vale $|a+b|=|a|+|b|$.
\begin{solution}
\textbf{Falsa.} Con $a=2$, $b=-2$, il primo membro è $0$, il secondo $4$.
\end{solution}
\part Fra due numeri razionali distinti esiste sempre un altro numero razionale.
\begin{solution}
\textbf{Vera.} Se $a<b$, la loro media $(a+b)/2$ è razionale e strettamente compresa fra $a$ e $b$.
\end{solution}
\end{parts}
\end{esercizio}

\begin{esercizio}
\obiettivo{interpretare opposto, modulo e distanza}
\question Rispondi motivando.
\begin{parts}
\part Sia $a\in\Z$, $a\ne0$. Quando $-a$ è positivo? $-|a|$ può essere positivo? $|-a|$ può essere negativo?
\begin{solution}
$-a$ è positivo se $a<0$. Poiché $|a|>0$, $-|a|<0$ e $|-a|=|a|>0$.
\end{solution}
\part Siano $a<0$, $b>0$ e $|a|<|b|$. Confronta $a$ con $-b$ e $-a$ con $b$. Esprimi con il modulo la distanza tra $-\frac32$ e $\frac54$ e calcolala.
\begin{solution}
$-a=|a|<|b|=b$, quindi $a>-b$. La distanza è $\left|\frac54-(-\frac32)\right|=\left|\frac{11}4\right|=\boxed{\frac{11}4}$.
\end{solution}
\end{parts}
\end{esercizio}

\begin{esercizio}
\obiettivo{confrontare numeri e valori assoluti}
\question Inserisci $<$, $=$ o $>$. Mostra i calcoli di M.C.D. e m.c.m. mediante scomposizione in fattori primi.
\begin{parts}
\part $\displaystyle |18-|-10||\ \ldots\ \operatorname{mcm}(2,8);\qquad \operatorname{MCD}(84,126)\ \ldots\ |-40|$.
\begin{solution}
$|18-10|=8=\operatorname{mcm}(2,8)$. $84=2^2\cdot3\cdot7$, $126=2\cdot3^2\cdot7$, quindi $\operatorname{MCD}=42>40$.
\end{solution}
\part $\displaystyle -|-7|\ \ldots\ -6;\qquad \operatorname{mcm}(12,18)\ \ldots\ |-6|^2;\qquad -3^2\ \ldots\ (-3)^2$.
\begin{solution}
$-7<-6$; $12=2^2\cdot3$ e $18=2\cdot3^2$, quindi il m.c.m. è $2^2\cdot3^2=36=6^2$; infine $-9<9$.
\end{solution}
\end{parts}
\end{esercizio}

\newpage

\begin{esercizio}
\obiettivo{calcolare espressioni in interi e usare le proprietà delle potenze}
\question Calcola mostrando i passaggi e applicando, dove possibile, le proprietà delle potenze.
\begin{parts}
\part $\displaystyle \bigl[(-2)^7\cdot(-2)^3:(-2)^8-3^2\bigr]^2-2\cdot|-8+3|-(-3)^3$.
\begin{solution}
Il prodotto e il quoziente danno $(-2)^{7+3-8}=4$. Quindi $(4-9)^2-2\cdot5+27=25-10+27=\boxed{42}$.
\end{solution}
\part $\displaystyle \left\{(-3)^6:\bigl[(-3)^2\bigr]^2+(-1)^{25}\right\}\cdot(-2)^3-\bigl[|-12+5|-2^3\bigr]^2$.
\begin{solution}
La potenza di potenza vale $(-3)^4$, perciò il quoziente è $(-3)^2=9$. Segue $(9-1)(-8)-(7-8)^2=-64-1=\boxed{-65}$.
\end{solution}
\end{parts}
\end{esercizio}

\begin{esercizio}
\obiettivo{usare equivalenza e scrittura decimale}
\question Completa e mostra i passaggi.
\begin{parts}
\part $\displaystyle \frac{-18}{24}=\frac{\square}{4}=\frac{15}{\square};\qquad \frac7{12}=\frac{\square}{60}$. Enuncia la proprietà invariantiva.
\begin{solution}
I numeri mancanti sono $-3$, $-20$, $35$. Moltiplicando o dividendo numeratore e denominatore per lo stesso numero diverso da zero si ottiene una frazione equivalente.
\end{solution}
\part Riduci ai minimi termini $\frac{45}{120}$ e $\frac{42}{175}$ e scrivile in forma decimale, eseguendo le divisioni.
\begin{solution}
$\frac{45}{120}=\frac38=0{,}375$; $\frac{42}{175}=\frac6{25}=0{,}24$.
\end{solution}
\end{parts}
\end{esercizio}

\begin{esercizio}
\obiettivo{passare fra frazioni, decimali e percentuali}
\question Trasforma in frazioni ridotte ai minimi termini e ordina in senso crescente.
\begin{parts}
\part $\displaystyle -1{,}25;\quad -\frac76;\quad 0{,}\overline3;\quad 35\%;\quad \frac38$.
\begin{solution}
$-\frac54<-\frac76<\frac13<\frac7{20}<\frac38$.\par Con denominatore $120$ i numeratori sono $-150,-140,40,42,45$.
\end{solution}
\part $\displaystyle -0{,}6;\quad -65\%;\quad \frac15;\quad 0{,}25;\quad 30\%$.
\begin{solution}
$-\frac{13}{20}<-\frac35<\frac15<\frac14<\frac3{10}$. Le conversioni sono $-0{,}6=-\frac35$, $-65\%=-\frac{13}{20}$, $0{,}25=\frac14$, $30\%=\frac3{10}$.
\end{solution}
\end{parts}
\end{esercizio}

\begin{esercizio}
\obiettivo{calcolare espressioni con frazioni, potenze negative e frazioni complesse}
\question Calcola applicando le proprietà delle potenze dove possibile. Mostra i passaggi e riduci il risultato.
\begin{parts}
\part $\displaystyle \left\{\left(-\frac23\right)^{-2}\cdot\left(\frac49\right)^2:\left(-\frac23\right)^3+\frac{1-\frac56}{\frac34-\frac12}\right\}:\left[\left(\frac32\right)^{-1}-\left(\frac12\right)^2\right]$.
\begin{solution}
Il prodotto e il quoziente valgono $\frac94\cdot\frac{16}{81}:(-\frac8{27})=-\frac32$; la frazione complessa vale $\frac{1/6}{1/4}=\frac23$. Quindi la graffa è $-\frac56$ e la quadra è $\frac23-\frac14=\frac5{12}$. Risultato: $-\frac56:\frac5{12}=\boxed{-2}$.
\end{solution}
\part $\displaystyle \left[\left(-\frac32\right)^{-3}:\left(-\frac32\right)^{-1}+\frac{\frac56-\frac12}{1-\frac34}\right]^{-1}-\left[\left(\frac23\right)^{-2}-\left(\frac12\right)^{-1}\right]^2$.
\begin{solution}
Il quoziente di potenze è $(-\frac32)^{-2}=\frac49$; la frazione complessa vale $\frac{1/3}{1/4}=\frac43$. La prima quadra vale $\frac{16}9$ e il suo inverso è $\frac9{16}$. La seconda quadra è $\frac94-2=\frac14$ e il quadrato vale $\frac1{16}$. Risultato $\boxed{\frac12}$.
\end{solution}
\end{parts}
\end{esercizio}

\newpage

\begin{esercizio}
\obiettivo{tradurre istruzioni verbali}
\question Scrivi un'unica espressione che traduca le istruzioni e calcolane il valore.
\begin{parts}
\part Sottrai al quadrato della somma di $-3$ e $5$ il prodotto tra $-2$ e il valore assoluto di $-7$.
\begin{solution}
$(-3+5)^2-(-2)\cdot|-7|=4+14=\boxed{18}$.
\end{solution}
\part Dividi la differenza tra i tre quarti e i cinque sesti per l'opposto di un terzo; aggiungi il quadrato di meno un mezzo.
\begin{solution}
$(\frac34-\frac56):(-\frac13)+(-\frac12)^2=\frac14+\frac14=\boxed{\frac12}$.
\end{solution}
\end{parts}
\end{esercizio}

\begin{esercizio}
\obiettivo{modellizzare una successione di operazioni}
\question Scrivi un'unica espressione per ciascun problema e calcolane il valore.
\begin{parts}
\part Giulia ha $756$ figurine. Regala un settimo del numero iniziale a un'amica e poi altre $156$ figurine al fratello. Infine regala un terzo di quelle rimaste. Quante ne conserva?
\begin{solution}
$\left[756(1-\frac17)-156\right](1-\frac13)=(648-156)\frac23=\boxed{328}$.
\end{solution}
\part Un serbatoio contiene $180$ litri d'acqua. Si utilizza un quarto dell'acqua iniziale, si aggiungono $18$ litri e infine si utilizzano i due noni dell'acqua presente. Quanti litri rimangono?
\begin{solution}
$\left[180(1-\frac14)+18\right](1-\frac29)=153\cdot\frac79=\boxed{119\text{ L}}$.
\end{solution}
\end{parts}
\end{esercizio}

\begin{esercizio}
\obiettivo{distinguere oggetti e intervalli}
\question Scrivi un'espressione e rispondi con l'unità di misura.
\begin{parts}
\part In una gara di $110$ m ci sono $10$ ostacoli equidistanti. Il primo è a $13{,}72$ m dalla partenza, l'ultimo a $14{,}02$ m dal traguardo. Quanto distano due ostacoli consecutivi?
\begin{solution}
Tra $10$ ostacoli vi sono $9$ intervalli. $(110-13{,}72-14{,}02):9=\boxed{9{,}14\text{ m}}$.
\end{solution}
\part Su un percorso di $100$ m ci sono $10$ ostacoli. Il primo è a $13$ m dalla partenza; due ostacoli consecutivi distano $8{,}5$ m. Quanto misura il tratto dall'ultimo ostacolo al traguardo?
\begin{solution}
$100-13-(10-1)\cdot8{,}5=\boxed{10{,}5\text{ m}}$.
\end{solution}
\end{parts}
\end{esercizio}

\begin{esercizio}
\obiettivo{trovare coincidenze con tempi non interi}
\question Risolvi mostrando come riconduci i tempi a numeri interi e spiegando perché la coincidenza trovata è la prima dopo la partenza.
\begin{parts}
\part Due segnali si accendono insieme. Il primo si riaccende ogni $5{,}5$ s, il secondo ogni $12{,}1$ s. Dopo quanti secondi si riaccenderanno insieme per la prima volta?
\begin{solution}
In decimi di secondo i periodi sono $55=5\cdot11$ e $121=11^2$. Il m.c.m. è $605$ decimi, cioè $\boxed{60{,}5\text{ s}}$. È il più piccolo tempo positivo multiplo di entrambi i periodi.
\end{solution}
\part Alice compie $12$ giri di pista in $8$ minuti e Bob $10$ giri in $6$ minuti, ad andatura costante. Partono insieme dalla linea di partenza. Dopo quanto tempo torneranno entrambi per la prima volta sulla linea? Quanti giri avranno completato?
\begin{solution}
Un giro dura $\frac23$ min $=40$ s per Alice, $\frac35$ min $=36$ s per Bob. $\operatorname{mcm}(40,36)=360$ s $=\boxed{6\text{ min}}$: Alice compie $\boxed9$ giri, Bob $\boxed{10}$. Il m.c.m. dà la prima coincidenza positiva sulla linea.
\end{solution}
\end{parts}
\end{esercizio}

\newpage

\begin{esercizio}
\obiettivo{usare scale e rapporti}
\question Risolvi mostrando l'impostazione.
\begin{parts}
\part Su una carta in scala $1:25\,000$ un sentiero misura $7{,}2$ cm. Quanto è lungo nella realtà, in km? Quanto misura sulla carta un tratto reale di $750$ m?
\begin{solution}
$7{,}2\cdot25\,000=180\,000$ cm $=\boxed{1{,}8\text{ km}}$. $750$ m $=75\,000$ cm, quindi sulla carta $\boxed{3\text{ cm}}$.
\end{solution}
\part Si mescolano sciroppo e acqua nel rapporto $2:7$. Per preparare $1{,}8$ litri di bevanda, quanto sciroppo e quanta acqua servono? Supponi che i volumi si sommino.
\begin{solution}
Le parti sono $9$, ciascuna da $1{,}8:9=0{,}2$ L. Servono $\boxed{0{,}4\text{ L}}$ di sciroppo e $\boxed{1{,}4\text{ L}}$ d'acqua.
\end{solution}
\end{parts}
\end{esercizio}

\begin{esercizio}
\obiettivo{calcolare parte, percentuale e totale}
\question Risolvi mostrando i calcoli.
\begin{parts}
\part In una scuola con $240$ studenti il $35\%$ partecipa a un torneo. Quanti sono i partecipanti? Se $21$ di loro disputano la finale, quale percentuale dei partecipanti raggiunge la finale?
\begin{solution}
$240\cdot\frac{35}{100}=\boxed{84}$ partecipanti. $\frac{21}{84}\cdot100\%=\boxed{25\%}$.
\end{solution}
\part Il $30\%$ degli iscritti a un corso corrisponde a $18$ persone. Quanti sono gli iscritti? Quale percentuale degli iscritti rappresentano $27$ persone?
\begin{solution}
$18:\frac{30}{100}=\boxed{60}$ iscritti. $\frac{27}{60}\cdot100\%=\boxed{45\%}$.
\end{solution}
\end{parts}
\end{esercizio}

\begin{esercizio}
\obiettivo{gestire variazioni percentuali successive}
\question Scrivi un'espressione e rispondi.
\begin{parts}
\part Una felpa costa $80$ euro. Il prezzo aumenta del $20\%$, poi diminuisce del $20\%$ rispetto al prezzo aumentato. Qual è il prezzo finale? Qual è la variazione percentuale complessiva rispetto al prezzo iniziale?
\begin{solution}
$80(1+\frac{20}{100})(1-\frac{20}{100})=80\cdot1{,}2\cdot0{,}8=\boxed{76{,}80\text{ euro}}$. La diminuzione di $3{,}20$ euro è il $\boxed{4\%}$ del prezzo iniziale.
\end{solution}
\part Un articolo viene scontato prima del $20\%$ e poi di un ulteriore $10\%$ sul prezzo già scontato. Alla fine costa $54$ euro. Quanto costava inizialmente? Qual è lo sconto percentuale complessivo?
\begin{solution}
La frazione di prezzo rimasta è $0{,}8\cdot0{,}9=0{,}72$. Prezzo iniziale $54:0{,}72=\boxed{75\text{ euro}}$. Sconto complessivo $1-0{,}72=0{,}28$, cioè $\boxed{28\%}$.
\end{solution}
\end{parts}
\end{esercizio}

\end{esercizi}
\finefila
\end{document}

